\chapter{System Implementation}

\section{Technology Stack Overview}

The system adopts a modern full-stack architecture designed for scalability, type safety, and efficient handling of financial data. Each component in the stack was selected based on its maturity, ecosystem support, and suitability for real-time analytical workloads.

To improve traceability, the versions listed in this section are aligned with the current project configuration files, which are frontend/package.json, backend/requirements.txt, and backend/runtime.txt at the time of writing.

\subsection{Frontend Technologies}
\begin{table}[H]
\centering
\begin{tabular}{lll}
\toprule
\textbf{Technology} & \textbf{Version} & \textbf{Purpose} \\
\midrule
Next.js & 14.2.5 & React framework with App Router for SSR/SSG \\
TypeScript & 5.x & Static typing and improved maintainability \\
Tailwind CSS & 4.x & Utility-first responsive UI styling \\
Recharts & 3.6.0 & Financial data visualization \\
TanStack Query & 5.90.x & Server-state management and caching \\
Axios & 1.13.x & HTTP client for API communication \\
\bottomrule
\end{tabular}
\caption{Frontend Technology Stack}
\label{tab:frontend_stack}
\end{table}

\subsection{Backend Technologies}

\begin{table}[H]
\centering
\begin{tabular}{lll}
\toprule
\textbf{Technology} & \textbf{Version} & \textbf{Purpose} \\
\midrule
Python & \(3.11.7\) & Core backend language \\
FastAPI & \(\ge 0.104.0\) & Asynchronous API framework \\
SQLAlchemy & \(\ge 2.0.23\) & ORM for database abstraction \\
Pydantic & \(\ge 2.5.0\) & Data validation and serialization \\
Pandas & \(\ge 1.5.0\) & Financial and time-series computation \\
APScheduler & \(\ge 3.10.4\) & Scheduled background tasks (news refresh jobs) \\
bcrypt & \(\ge 4.0.1\) & Secure password hashing \\
python-jose & \(\ge 3.3.0\) & JWT token handling \\
\bottomrule
\end{tabular}
\caption{Backend Technology Stack}
\label{tab:backend_stack}
\end{table}

\subsection{Database and External Services}

\begin{table}[H]
\centering
\begin{tabular}{lll}
\toprule
\textbf{Service} & \textbf{Provider} & \textbf{Purpose} \\
\midrule
PostgreSQL & Supabase & Relational data storage (5.7M+ records) \\
OpenAI API & OpenAI & LLM-powered educational chatbot responses \\
Alpha Vantage & Alpha Vantage & News sentiment analysis \\
FRED & Federal Reserve Economic Data & Macroeconomic indicators \\
Yahoo Finance & yfinance & Historical OHLCV price data \\
\bottomrule
\end{tabular}
\caption{Database and External Services}
\label{tab:external_services}
\end{table}

\subsection{Deployment Infrastructure}

\begin{table}[H]
\centering
\begin{tabular}{lll}
\toprule
\textbf{Component} & \textbf{Platform} & \textbf{Configuration} \\
\midrule
Frontend & Vercel & Edge CDN, GitHub CI/CD \\
Backend & Render & Docker, Gunicorn + Uvicorn \\
Database & Supabase & AWS ap-southeast-1, pooling enabled \\
\bottomrule
\end{tabular}
\caption{Deployment Infrastructure}
\label{tab:deployment_stack}
\end{table}



\section{Data Pipeline Implementation}

The data pipeline forms the backbone of the platform and is responsible for aggregating, cleaning, validating, and storing financial data across more than 2,200 ETFs.

\begin{figure}[H]
    \centering
    \includegraphics[width=0.95\textwidth]{fig/data_pipeline_flowchart.png}
    \caption{Data Pipeline Flowchart}
    \label{fig:data_pipeline}
\end{figure}

\subsection{Data Collection}

ETF metadata, historical prices, and news sentiment are collected from multiple authoritative sources to ensure data completeness and accuracy.

\textbf{ETF Metadata Collection:}
ETF metadata is collected using automated scripts that scrape ETFdb for U.S.-listed ETFs and utilize Yahoo Finance for Singapore-listed ETFs. All fields are normalized into a unified schema, ensuring consistency across markets.

\textbf{Historical Price Collection:}
The initial 15-year OHLCV dataset was collected via Jupyter notebooks using the yfinance library, exported to CSV files, and bulk-loaded into the database. Ongoing incremental updates are handled by app/services/price\_service.py, which is invoked automatically by the APScheduler job at application startup and by a lightweight CLI script for manual or backfill runs.

\begin{lstlisting}[style=pythonstyle, caption={Batch OHLCV Refresh via yfinance (price\_service.py)}, basicstyle=\ttfamily\footnotesize]
def refresh_all_prices(db, tickers=None, default_lookback_days=7):
    """Incrementally refresh OHLCV via yfinance (batched, skip up-to-date)."""
    all_tickers = tickers or _get_all_tickers(db)
    last_dates = _get_last_price_dates(db, all_tickers)

    today = datetime.now().date()
    needs_update = [
        (t, (last_dates[t] + timedelta(days=1)) if t in last_dates
         else (today - timedelta(days=default_lookback_days)))
        for t in all_tickers
        if (t not in last_dates or last_dates[t] + timedelta(days=1) < today)
    ]

    if not needs_update: return
    min_start = min(s for _, s in needs_update)

    for i in range(0, len(needs_update), 50):
        batch = needs_update[i:i+50]
        raw = yf.download(" ".join(t for t,_ in batch),
                          start=min_start, end=today,
                          group_by="ticker", auto_adjust=True, progress=False)

        parsed = _parse_download(raw, [t for t,_ in batch])
        for t, rows in parsed.items():
            _upsert_batch(db, t, rows, existing_set)
\end{lstlisting}

\textbf{News Data Collection:}
Alpha Vantage’s \texttt{NEWS\_SENTIMENT} endpoint is queried with application-level throttling. The implemented limiter enforces both per-minute and per-day caps to prevent quota violations, and parsed sentiment/relevance fields are persisted for downstream analytics.

\subsection{Automated Refresh Strategy}

While the initial bulk load covers historical data, several datasets require periodic refresh to remain accurate. Four APScheduler jobs are designed to handle this:

\begin{itemize}
    \item \textbf{OHLCV prices} (weekdays 20:00 ET): price\_service.refresh\_all\_prices() fetches only the missing date window per ticker via yfinance batch download (50 tickers/call), covering all 2,200+ ETFs. A CLI wrapper supports ad-hoc backfill runs.
    \item \textbf{News sentiment} (daily 18:00 ET): news\_service.refresh\_news\_for\_tickers() batches the top 100 ETFs by AUM into groups of 5 and calls Alpha Vantage \texttt{NEWS\_SENTIMENT}, capped at 25 calls/day and 5 calls/min.
    \item \textbf{ETF metadata} (monthly): A scheduled job re-queries ETFdb and \texttt{yfinance} to capture newly listed ETFs and updates to AUM, expense ratios, and category classifications, ensuring the ETF catalogue stays current.
    \item \textbf{Macro indicators} (monthly): A scheduled job pulls the latest FRED series (e.g.\ CPI, GDP, Federal Funds Rate) to incorporate newly released data points as economic indicators are published.
\end{itemize}

\subsection{Data Processing and Validation}

Before storage, raw data undergoes multiple preprocessing stages:
\begin{itemize}
    \item {Normalization:} Conversion of percentage and currency strings to numeric formats.
    \item {Gap Filling:} Forward-filling missing dates for continuous time series.
    \item{Deduplication:} Enforcement of composite keys on (ticker, date).
    \item {Validation:} Sanity checks on prices, dates, and ticker formats.
\end{itemize}

\subsection{Database Population}

Bulk data ingestion is optimized using PostgreSQL’s COPY command.
\begin{lstlisting}[
    style=bashstyle,
    caption={Bulk Data Upload Script},
    basicstyle=\ttfamily\footnotesize
]
psql $DATABASE_URL -c "\COPY etf_prices FROM 'prices.csv' WITH CSV HEADER"
\end{lstlisting}

% =========================================================
\section{Backend Implementation}

The backend is implemented using FastAPI and follows a service-oriented, modular architecture.

\subsection{Project Structure}

\begin{table}[H]
\centering
\begin{tabular}{p{4cm} p{10cm}}
\toprule
\textbf{File / Folder} & \textbf{Description} \\
\midrule
backend/ & Root folder \\
\hspace{0.5cm}app/ & Application package (FastAPI) \\
\hspace{1cm}main.py & App entry, scheduler, route registration \\
\hspace{1cm}auth.py & JWT, password hashing, auth dependencies \\
\hspace{1cm}database.py & SQLAlchemy engine, SessionLocal, Base, get\_db \\
\hspace{1cm}models.py & ORM models (etfs, etf\_prices, users, wallets, news) \\
\hspace{1cm}schemas.py & Pydantic request/response schemas \\
\hspace{1cm}routes/ & APIRouter modules (auth, etfs, prices, news, ...) \\
\hspace{1cm}services/ & Business logic and background helpers \\
\hspace{0.5cm}scripts/ & CLI utilities (e.g., update\_prices\_simple.py) \\
requirements.txt & Python dependencies \\
\bottomrule
\centering
\end{tabular}
\caption{Backend Project Structure}
\end{table}

The backend exposes a RESTful API organized into logical modules as defined in the system design (see Chapter 3, Table~\ref{tab:api-endpoints}). The implementation uses `APIRouter` to modularize the codebase, with specific routers for authentication, ETF data, portfolio management, and analytical services.

\subsection{Database Design}

The schema of the database is optimised for financial time-series data and rapid retrieval. Key tables include etf\_prices, which stores millions of records, partitioned and indexed for performance.

\begin{lstlisting}[language=SQL, caption={Core Tables Schema}, basicstyle=\ttfamily\footnotesize]
CREATE TABLE etf_prices (
    ticker VARCHAR(10) NOT NULL,
    date DATE NOT NULL,
    open DECIMAL(10, 4),
    high DECIMAL(10, 4),
    low DECIMAL(10, 4),
    close DECIMAL(10, 4),
    volume BIGINT,
    PRIMARY KEY (ticker, date)
);
CREATE INDEX idx_prices_ticker_date ON etf_prices (ticker, date DESC);

CREATE TABLE news_tickers (
    news_id INTEGER REFERENCES news(id),
    ticker VARCHAR(10) REFERENCES etfs(ticker),
    sentiment_score FLOAT,
    relevance_score FLOAT,
    PRIMARY KEY (news_id, ticker)
);
\end{lstlisting}

Other relational tables include etfs (ETF metadata and attributes such as AUM, expense ratio, category, and a snapshot of holdings), news and news\_tickers (news articles and their ticker-level sentiment/relevance links), macro\_indicators (FRED series values with timestamps), and user-related tables (users, portfolios, wallets, wallet\_holdings) which record user profiles, portfolio compositions, wallet goals, and holdings history. 


These tables are normalised where appropriate; referential integrity is enforced with foreign-key constraints (e.g., news\_tickers.news\_id -> news.id and wallet\_holdings.portfolio\_id -> portfolios.id). Indexes are created on frequently queried columns (for example, ticker, date, user\_id, portfolio\_id) to accelerate joins and lookups.


\subsection{Metrics Calculation}

Financial metrics are calculated on-demand using vectorized Pandas operations to ensure freshness and accuracy. These are some examples of functions that calculate financial metrics, not being collected via APIs, such as sharpe ratio and maximum drawdown.

\textbf{Sharpe Ratio:}
Measures risk-adjusted return using the 10-year Treasury yield as the risk-free rate.
\begin{lstlisting}[style=pythonstyle, caption={Sharpe Ratio Calculation}, basicstyle=\ttfamily\footnotesize]
def calculate_sharpe(returns: pd.Series, risk_free_rate: float = 0.04) -> float:
    annualized_return = returns.mean() * 252
    annualized_vol = returns.std() * math.sqrt(252)
    return (annualized_return - risk_free_rate) / annualized_vol
\end{lstlisting}

\textbf{Maximum Drawdown:}
Computes the largest peak-to-trough decline in the price history.
\begin{lstlisting}[style=pythonstyle, caption={Maximum Drawdown Calculation}, basicstyle=\ttfamily\footnotesize]
def calculate_max_drawdown(prices: pd.Series) -> float:
    running_max = prices.cummax()
    drawdown = (prices - running_max) / running_max
    return drawdown.min() * 100  # Return as percentage
\end{lstlisting}

\subsection{Risk Profiling}

User risk tolerance is derived from a questionnaire-driven scoring workflow implemented through backend quiz endpoints and frontend questionnaire pages. The current quiz contains 8 multiple-choice questions, with per-option points defined in "data/user\_info/quiz\_qus.json". Scores are summed and mapped into three operational profiles used by downstream recommendation logic: \textbf{conservative} (0--260), \textbf{balanced} (261--500), and \textbf{aggressive} (501--720). Because thresholds and point weights are loaded from configuration at runtime, profile boundaries can be adjusted without changing API contracts. When a quiz is submitted by an authenticated user, the latest profile is persisted to "users.risk\_profile" for reuse across later sessions.

\begin{table}[H]
\centering
\begin{tabular}{ll}
\toprule
\textbf{Profile} & \textbf{Practical Interpretation} \\
\midrule
Conservative & Capital preservation and lower volatility preference \\
Balanced & Mixed growth and stability preference \\
Aggressive & Higher return target with higher risk tolerance \\
\bottomrule
\end{tabular}
\caption{Implemented Risk Profile Categories}
\label{tab:risk_profiles}
\end{table}

More details about risk questions in Appendix~\ref{appendix:risk-questionnaire}.


\subsection{Recommendation Engine}

The current recommendation engine is implemented as a \textbf{rule-based multi-factor scorer} with profile-aware filtering in "routes/recommendations.py". It first removes ETFs in profile-specific avoided categories, then applies profile-specific beta hard filters where defined (for example conservative: beta \(\le\) 0.8, aggressive: beta \(\ge\) 1.0). Balanced recommendations are primarily shaped by scoring rather than strict beta exclusion.

Each candidate starts from a baseline score of 50 and receives additive bonuses or penalties based on parsed ETF attributes. The API returns ETFs with score \(\ge\) 40, sorted descending by score, and also supports optional category and region filters at query time.

The following listing is a simplified representation of the implemented scoring logic in the backend.

\textbf{Implemented scoring signals include:}
\begin{itemize}
    \item beta alignment with profile expectations,
    \item return and dividend heuristics (profile-dependent weighting),
    \item low expense ratio bonus (e.g., expense ratio < 0.2\%),
    \item category and asset-class matching bonuses,
    \item risk-category penalties for avoided categories.
\end{itemize}

\begin{lstlisting}[style=pythonstyle, caption={Rule-Based Recommendation Scoring (Simplified), basicstyle=\ttfamily\footnotesize}]
def calculate_recommendation_score(etf, profile):
    score = 50
    reasons = []

    beta = parse_beta(etf.beta)
    ytd_return = parse_percentage(etf.return_ytd)
    dividend = parse_percentage(etf.annual_dividend_yield)
    fee = parse_percentage(etf.expense_ratio)

    # Beta adjustments by profile
    if profile == "conservative":
        if beta is not None and beta < 0.5:
            score += 20; reasons.append("Very low volatility")
        elif beta is not None and beta < 0.8:
            score += 15; reasons.append("Low volatility")
    elif profile == "balanced":
        if beta is not None and 0.9 <= beta <= 1.1:
            score += 15; reasons.append("Market-aligned volatility")
    elif profile == "aggressive":
        if beta is not None and beta > 1.3:
            score += 20; reasons.append("High growth potential")
        elif beta is not None and beta > 1.1:
            score += 10; reasons.append("Above-market volatility")

    # Return/dividend heuristics
    if profile == "conservative" and dividend is not None and dividend > 3:
        score += 15
    if profile == "balanced" and dividend is not None and dividend > 2:
        score += 10
    if profile == "aggressive" and ytd_return is not None and ytd_return > 20:
        score += 20
    if profile == "balanced" and ytd_return is not None and ytd_return > 10:
        score += 10

    # Cost/category/class effects
    if fee is not None and fee < 0.2:
        score += 10
    if matches_preferred_category(etf, profile):
        score += 15
    if matches_preferred_asset_class(etf, profile):
        score += 10
    if in_avoid_category(etf, profile):
        score -= 30

    score = min(100, max(0, score))
    return score, reasons[:3]
\end{lstlisting}

\subsection{Scenario Analysis and VaR/CVaR}

The scenario module supports historical replay and stress tests via /scenarios/analyze at the portfolio, wallet, or single-ETF level. Implemented scenarios include a historical window (COVID-19 crash, 2020-02-19 to 2020-03-23) and uniform shocks (-10\%, -20\%, -30\%).

For each holding, the backend computes a base valuation (latest close or purchase/average cost) and reports scenario value, change (amount\%), and portfolio-level aggregate impact. The COVID replay also reports per-holding max drawdown.

The VaR endpoint (/scenarios/var) calculates VaR and CVaR using historical (empirical tail) and parametric (normal approximation) methods. Supported confidence levels are 0.90, 0.95, 0.99, with a configurable lookback window (30–1000 days).

\begin{lstlisting}[style=pythonstyle, caption={Implemented Portfolio VaR/CVaR Flow (Simplified)}, basicstyle=\ttfamily\footnotesize ]
def calculate_var(db, holdings, confidence_level=0.95, time_horizon_days=252):
    holding_values = {t: qty * latest_or_purchase_price(t) for t, qty in holdings}
    portfolio_value = sum(holding_values.values())
    portfolio_returns = aggregate_weighted_returns(holding_values, time_horizon_days)
    sorted_returns = sorted(portfolio_returns)

    var_cut = percentile(sorted_returns, 1 - confidence_level)
    hist_var_amt = abs(var_cut) * portfolio_value
    hist_tail = [r for r in sorted_returns if r <= var_cut]
    hist_cvar_amt = abs(mean(hist_tail)) * portfolio_value

    mu, sigma = mean(portfolio_returns), std(portfolio_returns)
    z = {0.90: 1.28, 0.95: 1.65, 0.99: 2.33}[confidence_level]
    para_var_amt = abs(mu - z * sigma) * portfolio_value
    para_cvar_amt = abs(mu - sigma * phi(z) / (1 - confidence_level)) * portfolio_value

    return build_response(hist_var_amt, para_var_amt, hist_cvar_amt, para_cvar_amt)
\end{lstlisting}

\subsection{AI Chatbot Service}

The chatbot integrates OpenAI LLM APIs for educational Q\&A. To reduce hallucinations and keep responses focused, the service follows retrieval-style grounding by injecting current ETF metadata and computed metrics into the prompt context before model invocation.

\begin{lstlisting}[style=pythonstyle, caption={AI System Prompt Construction}, basicstyle=\ttfamily\footnotesize]
def construct_prompt(etf_data, user_question):
    context = f"""
    You are a financial assistant analyzing {etf_data['name']} ({etf_data['ticker']}).
    Current Price: ${etf_data['price']} | Expense Ratio: {etf_data['expense_ratio']}%
    YTD Return: {etf_data['ytd_return']}% | Beta: {etf_data['beta']}
    Top Holdings: {', '.join(etf_data['top_holdings'])}
    """
    
    return [
        {"role": "system", "content": context + "\nProvide factual educational info. Do not give financial advice."},
        {"role": "user", "content": user_question}
    ]
\end{lstlisting}

\subsection{Authentication and Security}

Security is enforced using JWT (JSON Web Tokens) for stateless authentication and bcrypt for password hashing. This ensures that sensitive user credentials are never stored in plain text and that API routes are protected against unauthorized access.

\begin{lstlisting}[style=pythonstyle, caption={Password Hashing and Verification}, basicstyle=\ttfamily\footnotesize]
pwd_context = CryptContext(schemes=["bcrypt"], deprecated="auto")

def get_password_hash(password: str) -> str:
    return pwd_context.hash(password)

def verify_password(plain_password: str, hashed_password: str) -> bool:
    return pwd_context.verify(plain_password, hashed_password)
\end{lstlisting}

% =========================================================
\section{Frontend Implementation}

The frontend is developed using Next.js 14 with a component-driven architecture and responsive design principles.

\begin{longtable}{p{6cm} p{8cm}}
\toprule
\textbf{File / Folder} & \textbf{Description} \\
\midrule
\endfirsthead

\multicolumn{2}{c}%
{\tablename\ \thetable\ -- \textit{Continued from previous page}} \\
\toprule
\textbf{File / Folder} & \textbf{Description} \\
\midrule
\endhead

\midrule \multicolumn{2}{r}{\textit{Continued on next page}} \\
\endfoot

\bottomrule
\endlastfoot

frontend/ & Root folder \\
\hspace{0.5cm}src/ & Source code folder \\
\hspace{1cm}app/ & Next.js App Router pages \\
\hspace{1.5cm}page.tsx & Landing page (/) \\
\hspace{1.5cm}layout.tsx & Root layout with providers \\
\hspace{1.5cm}login/page.tsx & Login page \\
\hspace{1.5cm}signup/page.tsx & Signup page \\
\hspace{1.5cm}questionnaire/page.tsx & User questionnaire page \\
\hspace{1.5cm}recommendations/page.tsx & Recommendations page \\
\hspace{1.5cm}etfs/ & ETF discovery pages \\
\hspace{2cm}page.tsx & ETF discovery main page \\
\hspace{2cm}[ticker]/page.tsx & Individual ETF details page \\
\hspace{1.5cm}portfolio/page.tsx & Portfolio overview page \\
\hspace{1.5cm}wallets/page.tsx & Wallets overview page \\
\hspace{1.5cm}scenarios/page.tsx & Scenario analysis page \\
\hspace{1.5cm}profile/page.tsx & User profile page \\
\hspace{1cm}components/ & Reusable UI components \\
\hspace{1.5cm}AppShell.tsx & Main app shell layout \\
\hspace{1.5cm}Header.tsx & Header component \\
\hspace{1.5cm}NewsFeed.tsx & News feed component \\
\hspace{1.5cm}ScenarioAnalysis.tsx & Scenario analysis component \\
\hspace{1.5cm}VarAnalysis.tsx & VaR/CVaR analysis component \\
\hspace{1cm}contexts/ & React context providers \\
\hspace{1.5cm}AuthContext.tsx & Authentication context \\
\hspace{1cm}lib/ & Utility and helper libraries \\
\hspace{1.5cm}api.ts & Axios API client \\
\hspace{1.5cm}providers.tsx & Query/Auth providers \\
\hspace{0.5cm}package.json & NPM dependencies \\
\end{longtable}

\captionof{table}{Frontend Project Structure}

\subsection{Frontend Routes}

\begin{table}[H]
\centering
\label{tab:frontend_routes}
\begin{tabular}{lll}
\toprule
\textbf{Page} & \textbf{Route} & \textbf{Key Features} \\
\midrule
Landing & / & Feature overview and call-to-action \\
Login & /login & Authentication \\
Questionnaire & /questionnaire & Risk assessment wizard \\
Recommendations & /recommendations & Profile-aligned ETF ranking and rationale \\
ETF Discovery & /etfs & Search and filtering \\
ETF Details & /etfs/[ticker] & Charts, metrics, chatbot \\
Portfolio & /portfolio & Holdings, valuation, and news \\
Wallets & /wallets & Goal-based wallet management \\
Scenario Analysis & /scenarios & Stress testing and VaR/CVaR \\
Profile Settings & /profile & User profile and risk preference updates \\
\bottomrule
\end{tabular}
\caption{Frontend Route Map}
\end{table}


\section{Implementation Challenges}

Developing a data-intensive financial platform presented several significant technical challenges. This section details the key issues encountered and the solutions implemented.

\subsection{API Rate Limiting and Data Freshness}
\textbf{Challenge:} The Alpha Vantage API (free tier) restricts usage to 5 calls per minute and 500 per day. Updating news for 2,000+ ETFs daily was impossible under these constraints.

\textbf{Solution:} I implemented a prioritization queue using \texttt{APScheduler}. The system updates the top 100 ETFs by AUM daily (requiring 100 calls spread over 20 minutes) and updates other ETFs on-demand when a user views their details page. I also utilized caching strategies to prevent redundant API calls for the same ticker within a 24-hour window.

\subsection{Data Consistency and Gap Filling}
\textbf{Challenge:} Financial data from newly listed ETFs often contains null values or non-contiguous dates, breaking chart rendering and metric calculations (e.g., volatility requires continuous returns).

\textbf{Solution:} A robust preprocessing pipeline was built using Pandas. Missing prices on weekends/holidays are handled via forward-filling (\texttt{ffill()}), ensuring no look-ahead bias. For metrics requiring strict window sizes (e.g., 252-day volatility), ETFs with insufficient history are flagged and excluded from certain risk analytics to prevent skewed results.


\section{Deployment Configuration}

The platform is deployed using a cloud-native architecture with clearly separated frontend, backend, and database services to ensure scalability, reliability, and ease of maintenance.

\begin{itemize}
    \item {Backend:} FastAPI is containerized with Docker and deployed on Render using Gunicorn and Uvicorn workers.
    \item {Frontend:} The Next.js application is deployed on Vercel with automatic CI/CD from GitHub.
    \item {Database:} PostgreSQL on Supabase is configured with connection pooling, indexing, row-level security, and automated backups.
\end{itemize}




